% Deutsche Parallelfassung zu foundations/axioms.tex

\chapter{Axiome als Versagensgrenzen}
\label{chap:axiomatic-foundation}

Die begriffliche Struktur hat begründet, warum AIM ein Engineering Object von
seinen Zuständen, Beobachtungen, Bewertungen und Repräsentationen unterscheiden
muss. Die formale Aufgabe besteht nun darin, diese Unterscheidungen nutzbar zu
machen. Betrachten wir erneut ein Entwurfs-Alignment, das als CAD-Kurve
vorliegt, und einen vermessenen Zustand, der als abgetastete Polyline vorliegt.
Die beiden Geometrien können rechnerisch verglichen werden; Bedeutung erhält
dieser Vergleich jedoch erst, nachdem fünf häufige Versagensarten
ausgeschlossen wurden.

\section{Was das formale Modell nicht zusammenfallen lassen darf}

\begin{axiom}[Objektkontinuität]
\label{ax:object-continuity}
Die Identität eines Engineering Objects wird nicht durch einen Wechsel von
Zustand, Beobachtung, Bewertung oder Repräsentation ersetzt.
\end{axiom}

CAD-Kurve und Vermessungspolyline können ein Alignment, verschiedene Zustände
davon oder sogar verschiedene Objekte betreffen. Geometrische Ähnlichkeit kann
diese Identitätsfrage nicht entscheiden; die Identität muss bereits durch die
ingenieurtechnischen Qualifikationen begründet sein, unter denen der Vergleich
erfolgt.

\begin{axiom}[Trennung von Konstruktion und Repräsentation]
\label{ax:construction-representation-separation}
Konstruktives Ingenieurwissen und eine daraus abgeleitete Repräsentation
nehmen verschiedene formale Rollen ein.
\end{axiom}

Eine abgetastete Polyline kann Darstellung, Austausch oder Berechnung
unterstützen. Ihre Stützpunkte werden dadurch weder zur konstruktiven Definition
des Alignments, noch erzeugt eine erneute Abtastung ein neues autoritatives
Objekt.

\begin{axiom}[Trennung von Realität und Beobachtung]
\label{ax:reality-observation-separation}
Ein physischer Zustand und eine Beobachtung dieses Zustands sind verschiedene
Eingaben des ingenieurtechnischen Schließens.
\end{axiom}

Vermessungspunkte sind quellengebundene Evidenz. Sie bilden nicht selbst die
physische Eisenbahn, und Unsicherheit oder Unvollständigkeit der Beobachtung
darf nicht stillschweigend auf die physische Realität übertragen werden.

\begin{axiom}[Trennung von Bewertung und Entscheidung]
\label{ax:evaluation-decision-separation}
Eine Bewertung erzeugt ein qualifiziertes Ergebnis; sie stellt keine
Ingenieurentscheidung dar.
\end{axiom}

Eine berechnete Abweichung kann Annahme, Instandhaltungsplanung oder
Neuplanung unterstützen. Kein Grenzwert, Residuum oder Solver-Ergebnis erhält
Entscheidungsautorität allein dadurch, dass es berechnet wurde.

\begin{axiom}[Trennung von intrinsisch und abgeleitet]
\label{ax:intrinsic-derived-separation}
Abgeleitete Geometrie ersetzt nicht die intrinsische Konstruktion, aus der sie
gewonnen wurde.
\end{axiom}

Koordinaten, Tangenten, Versätze und abgetastete Stützpunkte sind nützliche
Folgen. Die formale Konstruktion muss deshalb die Abhängigkeiten offenlegen,
die sie erzeugen.

\section{Die konstruktive Abhängigkeit}

Es sei \(I\subseteq\mathbb{R}\) ein durch die Bogenlänge \(s\)
parametrisiertes Intervall. Ferner sei
\(\kappa:I\rightarrow\mathbb{R}\) die Krümmung; Richtungswinkel und ebene
Position werden mit \(\theta:I\rightarrow\mathbb{R}\) beziehungsweise
\(\gamma:I\rightarrow\mathbb{R}^{2}\) bezeichnet. Die folgenden
mathematischen Axiome formulieren die Abhängigkeit, durch die intrinsische
Konstruktion zu abgeleiteter Geometrie wird.

\begin{axiom}[Bogenlängenparametrisierung]
\label{ax:arc-length-parameterization}
\[
\lVert\gamma'(s)\rVert=1
\qquad\text{für alle }s\in I.
\]
\end{axiom}

\begin{axiom}[Entwicklung des Richtungswinkels]
\label{ax:heading-evolution}
\[
\theta'(s)=\kappa(s).
\]
\end{axiom}

\begin{axiom}[Tangentenrekonstruktion]
\label{ax:tangent-reconstruction}
\[
\gamma'(s)=
\begin{pmatrix}
\cos\theta(s)\\
\sin\theta(s)
\end{pmatrix}.
\]
\end{axiom}

\begin{proposition}[Krümmungsgetriebene Rekonstruktion]
\label{prop:curvature-driven-reconstruction}
Bei gegebener Anfangsposition \(\gamma(s_0)\), gegebenem Anfangswinkel
\(\theta(s_0)\) und einer hinreichend regulären Krümmungsfunktion \(\kappa\)
ist die ebene Trajektorie auf \(I\) eindeutig bestimmt.
\end{proposition}

\begin{proof}
Die Integration von \cref{ax:heading-evolution} ergibt
\[
\theta(s)=\theta(s_0)+\int_{s_0}^{s}\kappa(\sigma)\,\dd\sigma.
\]
Einsetzen in \cref{ax:tangent-reconstruction} und eine zweite Integration
ergeben
\[
\gamma(s)=\gamma(s_0)+
\int_{s_0}^{s}
\begin{pmatrix}
\cos\theta(\sigma)\\
\sin\theta(\sigma)
\end{pmatrix}
\dd\sigma.
\]
Die Anfangsbedingungen bestimmen die verbleibenden Integrationskonstanten.
\end{proof}

Dieses Ergebnis ist bewusst begrenzt. Es rekonstruiert eine ebene Trajektorie
aus benannten Eingaben und Annahmen. Es begründet keine Objektidentität,
beweist nicht, dass eine Vermessung denselben Zustand beobachtet, überträgt
keine Autorität auf eine Repräsentation und entscheidet nicht, ob eine
Abweichung akzeptabel ist.
